\section{Affine morphisms}
\label{section-affine}

\begin{definition}
\label{definition-affine}
A morphism of schemes $f : X \to S$ is called {\it affine} if
the inverse image of every affine open of $S$ is an affine
open of $X$.
\end{definition}

\begin{lemma}
\label{lemma-affine-separated}
An affine morphism is separated and quasi-compact.
\end{lemma}

\begin{proof}
Let $f : X \to S$ be affine. Quasi-compactness is immediate from
Schemes, Lemma \ref{schemes-lemma-quasi-compact-affine}.
We will show $f$ is separated using
Schemes, Lemma \ref{schemes-lemma-characterize-separated}. Let
$x_1, x_2 \in X$ be points of $X$ which map to the same point $s \in S$.
Choose any affine open $W \subset S$ containing $s$. By assumption
$f^{-1}(W)$ is affine. Apply the lemma cited with $U = V = f^{-1}(W)$.
\end{proof}

\begin{lemma}
\label{lemma-characterize-affine}
\begin{reference}
\cite[II, Corollary 1.3.2]{EGA}
\end{reference}
Let $f : X \to S$ be a morphism of schemes.
The following are equivalent
\begin{enumerate}
\item The morphism $f$ is affine.
\item There exists an affine open covering $S = \bigcup W_j$
such that each $f^{-1}(W_j)$ is affine.
\item There exists a quasi-coherent sheaf of $\mathcal{O}_S$-algebras
$\mathcal{A}$ and an isomorphism
$X \cong \underline{\Spec}_S(\mathcal{A})$
of schemes over $S$. See
Constructions, Section \ref{constructions-section-spec} for notation.
\end{enumerate}
Moreover, in this case $X = \underline{\Spec}_S(f_*\mathcal{O}_X)$.
\end{lemma}

\begin{proof}
It is obvious that (1) implies (2).

\medskip\noindent
Assume $S = \bigcup_{j \in J} W_j$ is an affine open covering such that
each $f^{-1}(W_j)$ is affine. By
Schemes, Lemma \ref{schemes-lemma-quasi-compact-affine} we see
that $f$ is quasi-compact. By
Schemes, Lemma \ref{schemes-lemma-characterize-quasi-separated}
we see the morphism $f$ is quasi-separated. Hence by
Schemes, Lemma \ref{schemes-lemma-push-forward-quasi-coherent} the
sheaf $\mathcal{A} = f_*\mathcal{O}_X$ is a quasi-coherent sheaf
of $\mathcal{O}_S$-algebras. Thus we have the scheme
$g : Y = \underline{\Spec}_S(\mathcal{A}) \to S$ over $S$.
The identity map
$\text{id} : \mathcal{A} = f_*\mathcal{O}_X \to f_*\mathcal{O}_X$
provides, via the definition of the relative spectrum,
a morphism $can : X \to Y$ over $S$, see
Constructions, Lemma \ref{constructions-lemma-canonical-morphism}.
By assumption and the lemma just cited
the restriction $can|_{f^{-1}(W_j)} : f^{-1}(W_j) \to g^{-1}(W_j)$
is an isomorphism. Thus $can$ is an isomorphism.
We have shown that (2) implies (3).

\medskip\noindent
Assume (3). By Constructions, Lemma \ref{constructions-lemma-spec-properties}
we see that the inverse image of every affine open is affine, and hence
the morphism is affine by definition.
\end{proof}

\begin{remark}
\label{remark-direct-argument}
We can also argue directly that (2) implies (1) in
Lemma \ref{lemma-characterize-affine} above as follows.
Assume $S = \bigcup W_j$ is an affine open covering
such that each $f^{-1}(W_j)$ is affine.
First argue that $\mathcal{A} = f_*\mathcal{O}_X$ is quasi-coherent
as in the proof above.
Let $\Spec(R) = V \subset S$ be affine open.
We have to show that $f^{-1}(V)$ is affine. Set
$A = \mathcal{A}(V) = f_*\mathcal{O}_X(V) = \mathcal{O}_X(f^{-1}(V))$.
By Schemes, Lemma \ref{schemes-lemma-morphism-into-affine} there is
a canonical morphism $\psi : f^{-1}(V) \to \Spec(A)$ over
$\Spec(R) = V$.
By Schemes, Lemma \ref{schemes-lemma-good-subcover} there exists
an integer $n \geq 0$, a standard open covering
$V = \bigcup_{i = 1, \ldots, n} D(h_i)$, $h_i \in R$, and a map
$a : \{1, \ldots, n\} \to J$ such that each $D(h_i)$ is also
a standard open of the affine scheme $W_{a(i)}$. The inverse image
of a standard open under a morphism of affine schemes is standard open, see
Algebra, Lemma \ref{algebra-lemma-spec-functorial}. Hence we see
that $f^{-1}(D(h_i))$ is a standard open of $f^{-1}(W_{a(i)})$,
in particular that $f^{-1}(D(h_i))$ is affine. Because $\mathcal{A}$
is quasi-coherent we have
$A_{h_i} = \mathcal{A}(D(h_i)) = \mathcal{O}_X(f^{-1}(D(h_i)))$,
so $f^{-1}(D(h_i))$ is the spectrum of $A_{h_i}$.
It follows that the morphism $\psi$ induces an isomorphism of the open
$f^{-1}(D(h_i))$ with the open $\Spec(A_{h_i})$ of
$\Spec(A)$. Since $f^{-1}(V) = \bigcup f^{-1}(D(h_i))$
and $\Spec(A) = \bigcup \Spec(A_{h_i})$ we win.
\end{remark}

\begin{lemma}
\label{lemma-affine-equivalence-algebras}
Let $S$ be a scheme. There is an anti-equivalence of categories
$$
\begin{matrix}
\text{Schemes affine} \\
\text{over }S
\end{matrix}
\longleftrightarrow
\begin{matrix}
\text{quasi-coherent sheaves} \\
\text{of }\mathcal{O}_S\text{-algebras}
\end{matrix}
$$
which associates to $f : X \to S$ the sheaf $f_*\mathcal{O}_X$.
Moreover, this equivalence is compatible with arbitrary base change.
\end{lemma}

\begin{proof}
The functor from right to left is given by $\underline{\Spec}_S$.
The two functors are mutually inverse by
Lemma \ref{lemma-characterize-affine} and
Constructions, Lemma \ref{constructions-lemma-spec-properties} part (3).
The final statement is
Constructions, Lemma \ref{constructions-lemma-spec-properties} part (2).
\end{proof}

\begin{lemma}
\label{lemma-quasi-coherence-scalar-restriction}
\begin{reference}
\cite[I, Proposition 9.6.1]{EGA}
\end{reference}
Let $S$ be a scheme and let $\mathcal{A}$ be a quasi-coherent
$\mathcal{O}_S$-algebra. An $\mathcal{A}$-module is
quasi-coherent as an $\mathcal{O}_S$-module if and only if
it is quasi-coherent as an $\mathcal{A}$-module.
\end{lemma}

\begin{proof}
Let $\mathcal{F}$ be an $\mathcal{A}$-module.
If $\mathcal{F}$ is quasi-coherent as an $\mathcal{A}$-module, then
for every $s \in S$ there exists an open neighbourhood $U$ of $s$
and an exact sequence
$$
\bigoplus\nolimits_J\mathcal{A}|_U\to
\bigoplus\nolimits_I\mathcal{A}|_U\to
\mathcal{F}|_U
$$
of $\mathcal{A}|_U$-modules.
Then this is also an exact sequence of $\mathcal{O}_U$-modules.
Hence $\mathcal{F}|_U$ is quasi-coherent as the cokernel of a morphism
of quasi-coherent $\mathcal{O}_U$-modules on a scheme. It follows
that $\mathcal{F}$ is quasi-coherent as an $\mathcal{O}_X$-module.

\medskip\noindent
Conversely, assume $\mathcal{F}$ is quasi-coherent as an $\mathcal{O}_X$-module.
Pick an open affine $\Spec(R) = U \subset S$. We have isomorphisms of
$\mathcal{O}_U$-modules $\mathcal{A}|_U \cong \widetilde{A}$ and
$\mathcal{F}|_U \cong \widetilde{M}$, for some $R$-algebra $A$
and some $R$-module $M$. The $\mathcal{A}$-module structure on $\mathcal{F}$
translates into an $A$-module structure on $M$ compatible with the
given $R$-module structure (details omitted). Choose an exact sequence
$$
\bigoplus\nolimits_J A \to
\bigoplus\nolimits_I A \to
M \to 0
$$
of $A$-modules. Since the functor $\widetilde{\ }$ is exact, this
produces an exact sequence
$$
\bigoplus\nolimits_J \widetilde{A}\to
\bigoplus\nolimits_I \widetilde{A}\to
\widetilde{M} \to 0
$$
of $\widetilde{A}$-modules. This means that
$\mathcal{F}$ is quasi-coherent as an $\mathcal{A}$-module.
\end{proof}

\begin{lemma}
\label{lemma-affine-equivalence-modules}
Let $f : X \to S$ be an affine morphism of schemes.
Let $\mathcal{A} = f_*\mathcal{O}_X$.
The functor $\mathcal{F} \mapsto f_*\mathcal{F}$ induces
an equivalence of categories
$$
\left\{
\begin{matrix}
\text{category of quasi-coherent}\\
\mathcal{O}_X\text{-modules}
\end{matrix}
\right\}
\longrightarrow
\left\{
\begin{matrix}
\text{category of quasi-coherent}\\
\mathcal{A}\text{-modules}
\end{matrix}
\right\}
$$
Moreover, an $\mathcal{A}$-module is
quasi-coherent as an $\mathcal{O}_S$-module if and only if
it is quasi-coherent as an $\mathcal{A}$-module.
\end{lemma}

\begin{proof}
The final statement is Lemma \ref{lemma-quasi-coherence-scalar-restriction}.
By Lemma \ref{lemma-affine-separated} and
Schemes, Lemma \ref{schemes-lemma-push-forward-quasi-coherent}
the pushforward $f_*\mathcal{F}$ of a quasi-coherent $\mathcal{O}_X$-module
$\mathcal{F}$
is a quasi-coherent $\mathcal{O}_S$-module. Hence a functor as in
the statement of the lemma.

\medskip\noindent
We will construct an quasi-inverse $h$ to this functor. Let $\mathcal{G}$
be a quasi-coherent $\mathcal{A}$-module. Then we set
$$
h(\mathcal{G}) = f^*\mathcal{G} \otimes_{f^*\mathcal{A}} \mathcal{O}_X =
f^*\mathcal{G} \otimes_{f^*f_*\mathcal{O}_X} \mathcal{O}_X
$$
Elucidation: the pullback $f^*\mathcal{A} = f^*f_*\mathcal{O}_X$
is an $\mathcal{O}_X$-algebra, the adjunction map
$f^*f_*\mathcal{O}_X \to \mathcal{O}_X$ is an algebra homomorphism, and
the pullback $f^*\mathcal{G}$ is an $f^*\mathcal{A}$-module. Observe that
$h(\mathcal{G})$ is quasi-coherent as quasi-coherence is preserved by
pullbacks and change of rings. Observe that there is a functorial map
$$
h(f_*\mathcal{F}) =
f^*f_*\mathcal{F} \otimes_{f^*f_*\mathcal{O}_X} \mathcal{O}_X \to \mathcal{F}
$$
coming from the adjunction map $f^*f_*\mathcal{F} \to \mathcal{F}$ and
a functorial map
$$
\mathcal{G} \to f_*h(\mathcal{G}) \quad\text{adjoint to the map}\quad
f^*\mathcal{G} \to f^*\mathcal{G} \otimes_{f^*f_*\mathcal{O}_X} \mathcal{O}_X
$$
which sends a local section $s$ of $f^*\mathcal{F}$ to $s \otimes 1$.
To finish the proof it suffices to show that these maps are isomorphisms
for $\mathcal{F}$ and $\mathcal{G}$ as above. This may be checked on the
members of an affine covering, i.e., when $X$ and $S$ are affine.

\medskip\noindent
The key algebra observation which makes this work is the
following: Let $R \to A$ be a ring map. Let $N$ be an $A$-module. Then
$$
(N \otimes_R A) \otimes_{(A \otimes_R A)} A = N
$$
Namely, the left hand side of this equality is the effect of applying
$h$ to the quasi-coherent $\widetilde{A}$-module $\widetilde{N}$ on $\Spec(R)$.
We omit the details.
\end{proof}

\begin{lemma}
\label{lemma-composition-affine}
The composition of affine morphisms is affine.
\end{lemma}

\begin{proof}
Let $f : X \to Y$ and $g : Y \to Z$ be affine morphisms.
Let $U \subset Z$ be affine open. Then $g^{-1}(U)$ is affine
by assumption on $g$. Whereupon $f^{-1}(g^{-1}(U))$ is affine
by assumption on $f$. Hence $(g \circ f)^{-1}(U)$ is affine.
\end{proof}

\begin{lemma}
\label{lemma-base-change-affine}
The base change of an affine morphism is affine.
\end{lemma}

\begin{proof}
Let $f : X \to S$ be an affine morphism. Let $S' \to S$ be any morphism.
Denote $f' : X_{S'} = S' \times_S X \to S'$ the base change of $f$.
For every $s' \in S'$ there exists an open affine neighbourhood
$s' \in V \subset S'$ which maps into some open affine $U \subset S$.
By assumption $f^{-1}(U)$ is affine. By the material in
Schemes, Section \ref{schemes-section-fibre-products} we see
that $f^{-1}(U)_V = V \times_U f^{-1}(U)$ is affine and equal
to $(f')^{-1}(V)$. This proves that $S'$ has an open covering by
affines whose inverse image under $f'$ is affine. We conclude
by Lemma \ref{lemma-characterize-affine} above.
\end{proof}

\begin{lemma}
\label{lemma-closed-immersion-affine}
A closed immersion is affine.
\end{lemma}

\begin{proof}
The first indication of this is
Schemes, Lemma \ref{schemes-lemma-closed-immersion-affine-case}.
See Schemes, Lemma \ref{schemes-lemma-closed-subspace-scheme}
for a complete statement.
\end{proof}

\begin{lemma}
\label{lemma-affine-s-open}
Let $X$ be a scheme.
Let $\mathcal{L}$ be an invertible $\mathcal{O}_X$-module.
Let $s \in \Gamma(X, \mathcal{L})$.
The inclusion morphism $j : X_s \to X$ is affine.
\end{lemma}

\begin{proof}
This follows from Properties, Lemma \ref{properties-lemma-affine-cap-s-open}
and the definition.
\end{proof}

\begin{lemma}
\label{lemma-affine-permanence}
Suppose $g : X \to Y$ is a morphism of schemes over $S$.
\begin{enumerate}
\item If $X$ is affine over $S$ and $\Delta : Y \to Y \times_S Y$ is affine,
then $g$ is affine.
\item If $X$ is affine over $S$ and $Y$ is separated over $S$,
then $g$ is affine.
\item A morphism from an affine scheme to a scheme with affine
diagonal is affine.
\item A morphism from an affine scheme to a separated scheme is affine.
\end{enumerate}
\end{lemma}

\begin{proof}
Proof of (1). The base change $X \times_S Y \to Y$ is affine by
Lemma \ref{lemma-base-change-affine}.
The morphism $(1, g) : X \to X \times_S Y$ is the base change of
$Y \to Y \times_S Y$ by the morphism $X \times_S Y \to Y \times_S Y$.
Hence it is affine by
Lemma \ref{lemma-base-change-affine}.
The composition of affine morphisms is affine
(see Lemma \ref{lemma-composition-affine}) and (1) follows.
Part (2) follows from (1) as a closed immersion is affine
(see Lemma \ref{lemma-closed-immersion-affine}) and $Y/S$ separated
means $\Delta$ is a closed immersion. Parts (3) and (4) are special
cases of (1) and (2).
\end{proof}

\begin{lemma}
\label{lemma-morphism-affines-affine}
A morphism between affine schemes is affine.
\end{lemma}

\begin{proof}
Immediate from Lemma \ref{lemma-affine-permanence} with
$S = \Spec(\mathbf{Z})$. It also follows directly from the
equivalence of (1) and (2) in Lemma \ref{lemma-characterize-affine}.
\end{proof}

\begin{lemma}
\label{lemma-Artinian-affine}
Let $S$ be a scheme.
Let $A$ be an Artinian ring.
Any morphism $\Spec(A) \to S$ is affine.
\end{lemma}

\begin{proof}
Omitted.
\end{proof}

\begin{lemma}
\label{lemma-get-affine}
Let $j : Y \to X$ be an immersion of schemes.
Assume there exists an open $U \subset X$ with complement
$Z = X \setminus U$ such that
\begin{enumerate}
\item $U \to X$ is affine,
\item $j^{-1}(U) \to U$ is affine, and
\item $j(Y) \cap Z$ is closed.
\end{enumerate}
Then $j$ is affine. In particular, if $X$ is affine, so is $Y$.
\end{lemma}

\begin{proof}
By Schemes, Definition \ref{schemes-definition-immersion} there exists an
open subscheme $W \subset X$ such that $j$ factors as a closed immersion
$i : Y \to W$ followed by the inclusion morphism $W \to X$.
Since a closed immersion is affine
(Lemma \ref{lemma-closed-immersion-affine}),
we see that for every $x \in W$ there is an affine open
neighbourhood of $x$ in $X$ whose inverse image under $j$ is affine.
If $x \in U$, then the same thing is true by assumption (2).
Finally, assume $x \in Z$ and $x \not \in W$. Then $x \not \in j(Y) \cap Z$.
By assumption (3) we can find an affine open neighbourhood
$V \subset X$ of $x$ which does not meet $j(Y) \cap Z$.
Then $j^{-1}(V) = j^{-1}(V \cap U)$ which is affine
by assumptions (1) and (2). It follows that $j$ is affine by
Lemma \ref{lemma-characterize-affine}.
\end{proof}







